\documentclass[11pt]{article}
\usepackage[a4paper,margin=1in]{geometry}
\usepackage{amsmath,amssymb}
\usepackage{hyperref}

\title{Status Note: High-Rank CBAP/OAP Failure Near \(p=2\)}
\author{fa\_banach\_001 lane 7}
\date{June 18, 2026}

\begin{document}
\maketitle

\paragraph{Status.}
This is a \texttt{literature\_implied\_answer (partial high-rank subcase)}.
It records a later theorem of de Laat--de la Salle \cite{dLdlS2014} that
partially answers an endpoint question raised in de Laat--de la Salle
\cite{dLdlS2015}.

\paragraph{Source question.}
In the introduction of \cite{dLdlS2015}, after improving the known failure of
CBAP and OAP for noncommutative \(L^p(L(\Gamma))\) spaces associated with
higher-rank lattices from the range
\[
  [1,12/11)\cup(12,\infty]
  \quad\hbox{to}\quad
  [1,10/9)\cup(10,\infty],
\]
the authors ask whether this failure can be extended to all \(p\neq 2\).

\paragraph{Supporting theorem.}
The later paper \cite{dLdlS2014}, Theorem 1.1, proves that if \(n\ge 3\) and
\(r\ge 2n-3\), then
\[
  L^p(L(\mathrm{SL}(r,\mathbb Z)))
\]
does not have CBAP for
\[
  p\in [1,2-2/n)\cup(2+2/(n-2),\infty].
\]
The paper also notes that the same argument gives failure of OAP, and Theorem
1.3 gives the corresponding conclusion for lattices in connected simple real
Lie groups of sufficiently high real rank.

\paragraph{Identification.}
Fix \(p\neq 2\).  If \(p>2\), choose \(n\) so large that
\(p>2+2/(n-2)\).  If \(1\le p<2\), choose \(n\) so large that
\(p<2-2/n\).  Then Theorem 1.1 of \cite{dLdlS2014} applies to
\(\mathrm{SL}(r,\mathbb Z)\) for every \(r\ge 2n-3\), and Theorem 1.3 applies
to sufficiently high-rank real simple Lie groups and their lattices.  Thus the
source endpoint question has an affirmative answer in the rank-varying
high-rank subcase: for every \(p\neq 2\), high enough rank lattices have
noncommutative \(L^p\)-spaces without CBAP and without OAP.

\paragraph{Scope limitations.}
This does not solve the fixed low-rank problem.  In particular, \cite{dLdlS2014}
explicitly leaves open whether \(L^p(L(\mathrm{SL}(3,\mathbb Z)))\) has CBAP
for some \(p\neq 2\).  It also does not settle the separate question in
\cite{dLdlS2015} about whether the classes \(\mathcal E_r\) contain all Banach
spaces of nontrivial type.

\begin{thebibliography}{9}
\bibitem{dLdlS2015}
T. de Laat and M. de la Salle,
\emph{Strong property (T) for higher rank simple Lie groups},
Proc. Lond. Math. Soc. (3) 111 (2015), no. 4, 936--966;
arXiv:1401.3611.

\bibitem{dLdlS2014}
T. de Laat and M. de la Salle,
\emph{Approximation properties for noncommutative \(L^p\)-spaces of high rank
lattices and nonembeddability of expanders},
J. Reine Angew. Math. 737 (2018), 49--69;
arXiv:1403.6415.
\end{thebibliography}

\end{document}
